% Generated draft: requires mathematical and source review.
\documentclass[11pt]{article}
\usepackage[margin=25mm]{geometry}
\usepackage{fontspec}
\setmainfont{Latin Modern Roman}
\usepackage{amsmath,amssymb,mathtools,unicode-math}
\setmathfont{Latin Modern Math}
\usepackage{enumitem,needspace,xurl,hyperref}
\setlength{\parindent}{0pt}
\setlength{\parskip}{5pt}
\setlength{\emergencystretch}{3em}
\newcommand{\N}{\mathbb{N}}
\newcommand{\Z}{\mathbb{Z}}
\newcommand{\Q}{\mathbb{Q}}
\newcommand{\R}{\mathbb{R}}
\newcommand{\C}{\mathbb{C}}
\newcommand{\F}{\mathbb{F}}
\newcommand{\A}{\mathbb{A}}
\newcommand{\PP}{\mathbb{P}}
\newcommand{\E}{\mathbb{E}}
\newcommand{\eps}{\varepsilon}
\newcommand{\dd}{\,\mathrm{d}}
\begin{document}
\section*{UnsolvedMath ID: 30001415}
\textbf{Source problem code:} catalogue entry
\subsection{Definitions and mathematical statement}
For \(1<p<\infty\), the \(p\)-Laplacian is
\[
\Delta_p u=\operatorname{div}(|\nabla u|^{p-2}\nabla u).
\]
The supplied record asks for the structure of second eigenfunctions on bounded domains, especially their nodal sets. It does not specify boundary conditions or the convention for the second eigenvalue, so it does not determine a unique spectral problem.

An explicitly conditional homogeneous Dirichlet formulation is as follows. Let \(\Omega\subset\R^n\) be a bounded open connected set. Here \(W_0^{1,p}(\Omega)\) denotes the closure of \(C_c^\infty(\Omega)\) in the Sobolev space \(W^{1,p}(\Omega)\). A Dirichlet eigenpair consists of \(\lambda\in\R\) and \(0\ne u\in W_0^{1,p}(\Omega)\) such that
\[
\int_\Omega |\nabla u|^{p-2}\nabla u\cdot\nabla\varphi\,\dd x
=\lambda\int_\Omega |u|^{p-2}u\varphi\,\dd x
\quad\text{for every }\varphi\in W_0^{1,p}(\Omega).
\]
Under this interpretation, set
\[
\lambda_1(\Omega,p)=\inf_{0\ne v\in W_0^{1,p}(\Omega)}
\frac{\int_\Omega |\nabla v|^p\,\dd x}{\int_\Omega |v|^p\,\dd x},
\]
and interpret \(\lambda_2(\Omega,p)\) as the least Dirichlet eigenvalue strictly greater than \(\lambda_1(\Omega,p)\). A second eigenfunction is an eigenfunction for \(\lambda_2(\Omega,p)\); it may be normalized by \(\int_\Omega |u|^p\,\dd x=1\). For its continuous interior representative, define the interior nodal set by
\[
Z_\Omega(u)=\{x\in\Omega:u(x)=0\}.
\]
The boundary is excluded from this definition.

\textbf{Requested structural problem, conditional on the preceding spectral interpretation.} Determine the possible second eigenfunctions and their interior nodal sets on bounded domains. In particular, establish or refute the following geometric descriptions in dimension two:
\begin{itemize}
\item For a disk \(D_R=\{x\in\R^2:|x|<R\}\), with \(R>0\), and every \(p\in(1,\infty)\), the nodal line is a diameter. Interpreting this as equality of the entire interior nodal set means
\[
Z_{D_R}(u)=D_R\cap\{x:x\cdot e=0\}
\quad\text{for some }e\in\R^2\text{ with }|e|=1.
\]
\item For the square \(S_a=(-a,a)^2\), with \(a>0\), the nodal line is vertical or horizontal when \(1<p<2\), and diagonal when \(p>2\). The corresponding entire-nodal-set descriptions are
\[
Z_{S_a}(u)=S_a\cap\{x_1=c\}
\quad\text{or}\quad
Z_{S_a}(u)=S_a\cap\{x_2=c\},\qquad c\in(-a,a),
\]
for \(1<p<2\), and
\[
Z_{S_a}(u)=S_a\cap\{x_2=x_1\}
\quad\text{or}\quad
Z_{S_a}(u)=S_a\cap\{x_2=-x_1\}
\]
for \(p>2\). The source does not explicitly say whether the vertical or horizontal segment must be a central axis, which would require \(c=0\).
\end{itemize}
For each geometric description, the source leaves unresolved whether it asserts existence of a second eigenfunction with that nodal set or asserts that every second eigenfunction has a nodal set of that form. These are distinct formulations and neither is selected here. Equality with the entire interior zero set is also an explicit interpretation of the source's phrase ``nodal line,'' rather than a verified source qualification. No square claim at \(p=2\) is supplied. The Dirichlet boundary condition and spectral convention above are conditional explanations, not hypotheses established by the supplied record.
\subsection{Short English statement}
The problem asks how second \(p\)-Laplacian eigenfunctions vanish: along a diameter in a disk for every \(1<p<\infty\), and along a vertical or horizontal segment in a square for \(1<p<2\), versus a diagonal for \(p>2\). Boundary conditions, the meaning of ``second,'' and whether these claims concern some or all second eigenfunctions need clarification. The displayed Dirichlet formulation is conditional. The supplied status is partially solved; this draft does not independently verify that assessment.
\subsection{Sources}
\begin{itemize}
\item Supplied local UnsolvedMath registry record, permanent ID \(30001415\), source problem code \texttt{catalogue entry}, entitled ``Nodal Geometry of Second \(p\)-Laplacian Eigenfunctions,'' in the Oberwolfach Reports Open Problems collection. The supplied proposed year is 2010; no individual proposer is identified.
\item Supplied report citation: \emph{Optimal Constants in the Theory of Sobolev Spaces and PDEs} (2010), Oberwolfach Reports, pp.~325--371, DOI \texttt{10.4171/owr/2010/08}. \url{https://doi.org/10.4171/owr/2010/08}. The supplied background describes these reports as records of mathematical workshops and their open-problem discussions.
\item Supplied literature reference: Cuesta--de Figueiredo--Gossez, 1999. \url{https://doi.org/10.1006/JDEQ.1999.3645}. No article title or more specific result citation is supplied.
\item Supplied literature reference: Oberwolfach Report 8/2010. \url{https://doi.org/10.4171/owr/2010/08}.
\item The supplied literature assessment was checked on \texttt{2026-08-21} and labels the problem partially solved. It states that the second \(p\)-Laplacian eigenvalue has variational and Fu\v{c}\'{i}k-spectrum characterizations and sign-changing nodal information, while these results do not exhaust the report's broad structural question. This is supplied dated curation derived from the OWR record, not independently verified here; the record requests independent verification of time-sensitive status and cited claims.
\item The registry labels the statement \texttt{corrected\_verified} and reports a manual semantic, typographic, and LaTeX review of its cleanup. These are supplied editorial assertions, not independent mathematical or source verification by this draft.
\end{itemize}
\end{document}
